% GENERATED FILE. Edit canonical lessons, then run npm run generate:books.
% canonical-source-hash: fnv1a64:90765d8bddfbdec5
% canonical-lessons: GU-C98-vaisakh, GU-C98-jeth, GU-C98-sravan, GU-C98-bhadarvo, GU-C98-aso

\chapter{Vaishakh, Jeth, Shravan, Bhadarvo, Aso}
\label{ch:gu-vaishakh-jeth-shravan-bhadarvo-aso}
\hlchaptermodality{98}

\begin{chapteropening}
\textbf{By the end of this chapter:} \emph{I can say Vaishakh, Jeth, Shravan, Bhadarvo and Aso.}

Complete the last lesson of chapter 98: I can say Vaishakh, Jeth, Shravan, Bhadarvo and Aso.
\end{chapteropening}

\section[\texorpdfstring{\gu{વૈશાખ} (vaiśākh)}{વૈશાખ (vaiśākh)}]{\gu{વૈશાખ} (vaiśākh) — Vaishakh, a Gujarati month}

\label{lesson:GU-C98-vaisakh}



\begin{quote}
Before the new one: say the Gujarati for Phagan, then the Gujarati for Chaitra.
\end{quote}

\subsection*{You'll want to know: \gu{વૈશાખ}}
\textbf{\gu{વૈશાખ}} (\emph{vaiśākh}) — \textquotedblleft{}Vaishakh, a Gujarati month\textquotedblright{}.

\begin{cousinweb}
\textbf{\gu{વૈશાખ}} (\emph{vaiśākh}) is Sanskrit \emph{vaiśākha}, \textquotedblleft{}the month April-May\textquotedblright{}, taken over whole as a learned loan, a tatsama (CDIAL 12149).
\end{cousinweb}

\begin{grammarlens}[title={What you've built}]
One more word for this chapter.
\end{grammarlens}

\subsection*{Your turn}
\begin{itemize}
\raggedright
  \item \emph{Say it:} \emph{vaiśākh}
  \item \emph{Say it:} \emph{vaiśākh}, once more
  \item \emph{Say it:} say \emph{vaiśākh}
  \item \emph{Recall:} say the Gujarati for Phagan, then the Gujarati for Chaitra, then say \emph{vaiśākh} again
\end{itemize}

\begin{tcolorbox}[breakable,colback=teal!4,colframe=teal!35!black,title={Before you move on}]
What does \gu{વૈશાખ} mean? (\textquotedblleft{}Vaishakh, a Gujarati month\textquotedblright{}.) Say it once more.
\end{tcolorbox}
\section[\texorpdfstring{\gu{જેઠ} (jeṭh)}{જેઠ (jeṭh)}]{\gu{જેઠ} (jeṭh) — Jeth, a Gujarati month}

\label{lesson:GU-C98-jeth}



\begin{quote}
Before the new one: say the Gujarati for Chaitra, then the Gujarati for Vaishakh.
\end{quote}

\subsection*{You'll want to know: \gu{જેઠ}}
\textbf{\gu{જેઠ}} (\emph{jeṭh}) — \textquotedblleft{}Jeth, a Gujarati month\textquotedblright{}.

\begin{cousinweb}
\textbf{\gu{જેઠ}} (\emph{jeṭh}) is a tadbhava, a word inherited from Sanskrit and reshaped by sound change. It comes from \emph{jyaiṣṭha}, \textquotedblleft{}the month May-June\textquotedblright{} (CDIAL 5293). It is named after \emph{jyeṣṭhā}, a lunar mansion.
\end{cousinweb}

\begin{grammarlens}[title={What you've built}]
One more word for this chapter.
\end{grammarlens}

\subsection*{Your turn}
\begin{itemize}
\raggedright
  \item \emph{Say it:} \emph{jeṭh}
  \item \emph{Say it:} \emph{jeṭh}, once more
  \item \emph{Say it:} say \emph{jeṭh}
  \item \emph{Recall:} say the Gujarati for Chaitra, then the Gujarati for Vaishakh, then say \emph{jeṭh} again
\end{itemize}

\begin{tcolorbox}[breakable,colback=teal!4,colframe=teal!35!black,title={Before you move on}]
What does \gu{જેઠ} mean? (\textquotedblleft{}Jeth, a Gujarati month\textquotedblright{}.) Say it once more.
\end{tcolorbox}
\section[\texorpdfstring{\gu{શ્રાવણ} (śrāvaṇ)}{શ્રાવણ (śrāvaṇ)}]{\gu{શ્રાવણ} (śrāvaṇ) — Shravan, a Gujarati month}

\label{lesson:GU-C98-sravan}



\begin{quote}
Before the new one: say the Gujarati for Vaishakh, then the Gujarati for Jeth.
\end{quote}

\subsection*{You'll want to know: \gu{શ્રાવણ}}
\textbf{\gu{શ્રાવણ}} (\emph{śrāvaṇ}) — \textquotedblleft{}Shravan, a Gujarati month\textquotedblright{}.

\begin{cousinweb}
\textbf{\gu{શ્રાવણ}} (\emph{śrāvaṇ}) is Sanskrit \emph{śrāvaṇa}, \textquotedblleft{}the twentieth or twenty-third asterism\textquotedblright{}, taken over whole as a learned loan, a tatsama (CDIAL 12699). CDIAL lists the inherited Gujarati form as \emph{sāvaṇ}, \textquotedblleft{}the fourth month\textquotedblright{}.
\end{cousinweb}

\begin{grammarlens}[title={What you've built}]
One more word for this chapter.
\end{grammarlens}

\subsection*{Your turn}
\begin{itemize}
\raggedright
  \item \emph{Say it:} \emph{śrāvaṇ}
  \item \emph{Say it:} \emph{śrāvaṇ}, once more
  \item \emph{Say it:} say \emph{śrāvaṇ}
  \item \emph{Recall:} say the Gujarati for Vaishakh, then the Gujarati for Jeth, then say \emph{śrāvaṇ} again
\end{itemize}

\begin{tcolorbox}[breakable,colback=teal!4,colframe=teal!35!black,title={Before you move on}]
What does \gu{શ્રાવણ} mean? (\textquotedblleft{}Shravan, a Gujarati month\textquotedblright{}.) Say it once more.
\end{tcolorbox}
\section[\texorpdfstring{\gu{ભાદરવો} (bhādarvo)}{ભાદરવો (bhādarvo)}]{\gu{ભાદરવો} (bhādarvo) — Bhadarvo, a Gujarati month}

\label{lesson:GU-C98-bhadarvo}



\begin{quote}
Before the new one: say the Gujarati for Jeth, then the Gujarati for Shravan.
\end{quote}

\subsection*{You'll want to know: \gu{ભાદરવો}}
\textbf{\gu{ભાદરવો}} (\emph{bhādarvo}) — \textquotedblleft{}Bhadarvo, a Gujarati month\textquotedblright{}.

\begin{cousinweb}
\textbf{\gu{ભાદરવો}} (\emph{bhādarvo}) is a tadbhava, a word inherited from Sanskrit and reshaped by sound change. It comes from \emph{bhādrapada}, \textquotedblleft{}the fifth month, August-September\textquotedblright{} (CDIAL 9447).
\end{cousinweb}

\begin{grammarlens}[title={What you've built}]
One more word for this chapter.
\end{grammarlens}

\subsection*{Your turn}
\begin{itemize}
\raggedright
  \item \emph{Say it:} \emph{bhādarvo}
  \item \emph{Say it:} \emph{bhādarvo}, once more
  \item \emph{Say it:} say \emph{bhādarvo}
  \item \emph{Recall:} say the Gujarati for Jeth, then the Gujarati for Shravan, then say \emph{bhādarvo} again
\end{itemize}

\begin{tcolorbox}[breakable,colback=teal!4,colframe=teal!35!black,title={Before you move on}]
What does \gu{ભાદરવો} mean? (\textquotedblleft{}Bhadarvo, a Gujarati month\textquotedblright{}.) Say it once more.
\end{tcolorbox}
\section[\texorpdfstring{\gu{આસો} (āso)}{આસો (āso)}]{\gu{આસો} (āso) — Aso, the month of Diwali}

\label{lesson:GU-C98-aso}



\begin{quote}
Before the new one: say the Gujarati for Shravan, then the Gujarati for Bhadarvo.
\end{quote}

\subsection*{You'll want to know: \gu{આસો}}
\textbf{\gu{આસો}} (\emph{āso}) — \textquotedblleft{}Aso, the month of Diwali\textquotedblright{}.

\begin{cousinweb}
\textbf{\gu{આસો}} (\emph{āso}) is a tadbhava of Sanskrit \emph{aśvayuja}, \textquotedblleft{}the month Āśvina, September-October\textquotedblright{} (CDIAL 925). CDIAL compares \emph{aśvayuj}, \textquotedblleft{}harnessing horses\textquotedblright{}, also the name of a constellation.
\end{cousinweb}

\begin{grammarlens}[title={What you've built}]
One more word for this chapter.
\end{grammarlens}

\subsection*{Your turn}
\begin{itemize}
\raggedright
  \item \emph{Say it:} \emph{āso}
  \item \emph{Say it:} \emph{āso}, once more
  \item \emph{Say it:} say \emph{āso}
  \item \emph{Recall:} say the Gujarati for Shravan, then the Gujarati for Bhadarvo, then say \emph{āso} again
\end{itemize}

\begin{tcolorbox}[breakable,colback=teal!4,colframe=teal!35!black,title={Before you move on}]
What does \gu{આસો} mean? (\textquotedblleft{}Aso, the month of Diwali\textquotedblright{}.) Say it once more.
\end{tcolorbox}
